\chapter{Relativistische Kinematik}

\section{Ziel dieses Kapitels}

Dieses Kapitel behandelt die relativistische Kinematik, die in der Kern- und
Teilchenphysik benoetigt wird.

Bei hohen Energien reicht die klassische Naeherung
\[
    E_{\mathrm{kin}}=\frac{1}{2}mv^2
\]
nicht mehr aus.
Stattdessen muessen Energie und Impuls relativistisch behandelt werden.

\begin{notebox}
Dieses Kapitel behandelt die Stichwoerter:
\begin{itemize}
    \item relativistische Energie,
    \item relativistischer Impuls,
    \item Viererimpuls,
    \item invariante Masse,
    \item Schwerpunktsenergie,
    \item Schwerpunktssystem,
    \item Laborsystem,
    \item Fixed-Target-Experiment,
    \item Collider,
    \item Schwellenenergie,
    \item Resonanzen,
    \item Luminositaet,
    \item Ereignisrate.
\end{itemize}
\end{notebox}

\section{Warum relativistische Kinematik?}

In Teilchenbeschleunigern erreichen Teilchen Energien, die vergleichbar mit oder
viel groesser als ihre Ruheenergie sind.
Dann bewegen sie sich mit Geschwindigkeiten nahe der Lichtgeschwindigkeit.

\begin{conceptbox}
Bei hohen Energien wird die Geschwindigkeit eines massiven Teilchens nicht beliebig
viel groesser.
Sie naehert sich der Lichtgeschwindigkeit $c$ an.

Zusaetzliche Energie erhoeht dann vor allem:
\begin{itemize}
    \item den Impuls,
    \item die Gesamtenergie,
    \item die im Schwerpunktsystem verfuegbare Energie.
\end{itemize}
\end{conceptbox}

\begin{examtipbox}
Bei relativistischen Teilchen gilt:

Mehr Energie bedeutet nicht wesentlich mehr Geschwindigkeit.

Mehr Energie bedeutet vor allem mehr Impuls und mehr moegliche Energie fuer neue
Teilchen im Schwerpunktsystem.
\end{examtipbox}

\section{Lorentzfaktor}

\begin{definitionbox}
Der \emph{Lorentzfaktor} ist definiert als
\[
    \gamma
    =
    \frac{1}{\sqrt{1-\beta^2}},
\]
wobei
\[
    \beta=\frac{v}{c}
\]
ist.
\end{definitionbox}

\begin{conceptbox}
Wenn
\[
    v \ll c
\]
ist, dann gilt
\[
    \beta \ll 1
\]
und
\[
    \gamma \approx 1.
\]
Dann ist die klassische Mechanik oft eine gute Naeherung.

Wenn
\[
    v \approx c
\]
ist, wird
\[
    \gamma
\]
gross.
Dann muss man relativistisch rechnen.
\end{conceptbox}

\begin{formulabox}
\[
    \beta=\frac{v}{c},
    \qquad
    \gamma=\frac{1}{\sqrt{1-\beta^2}}.
\]
\end{formulabox}

\section{Relativistische Energie}

\begin{definitionbox}
Die \emph{Ruheenergie} eines Teilchens mit Masse $m$ ist
\[
    E_0=mc^2.
\]
\end{definitionbox}

\begin{definitionbox}
Die \emph{relativistische Gesamtenergie} ist
\[
    E=\gamma mc^2.
\]
\end{definitionbox}

\begin{formulabox}
Die kinetische Energie ist die Differenz zwischen Gesamtenergie und Ruheenergie:
\[
    E_{\mathrm{kin}}
    =
    E-E_0
    =
    (\gamma-1)mc^2.
\]
\end{formulabox}

\begin{conceptbox}
Die Gesamtenergie besteht aus Ruheenergie und kinetischer Energie:
\[
    E = mc^2 + E_{\mathrm{kin}}.
\]

Bei kleinen Geschwindigkeiten geht die relativistische Formel in die klassische
Formel ueber:
\[
    E_{\mathrm{kin}}\approx \frac{1}{2}mv^2.
\]
\end{conceptbox}

\section{Relativistischer Impuls}

\begin{definitionbox}
Der relativistische Impuls eines Teilchens ist
\[
    \vec p = \gamma m \vec v.
\]
\end{definitionbox}

\begin{conceptbox}
Bei kleinen Geschwindigkeiten gilt:
\[
    \gamma\approx 1.
\]
Dann wird
\[
    \vec p\approx m\vec v.
\]
Das ist der klassische Impuls.

Bei relativistischen Geschwindigkeiten wird der Impuls deutlich groesser als
klassisch erwartet.
\end{conceptbox}

\begin{formulabox}
\[
    \vec p = \gamma m\vec v.
\]
\end{formulabox}

\section{Energie-Impuls-Beziehung}

Die wichtigste relativistische Beziehung zwischen Energie, Impuls und Masse lautet:

\begin{formulabox}
\[
    E^2
    =
    p^2c^2 + m^2c^4.
\]
\end{formulabox}

\begin{conceptbox}
Diese Gleichung gilt fuer freie Teilchen.

Spezialfaelle:
\begin{itemize}
    \item Teilchen in Ruhe:
    \[
        p=0
        \quad\Rightarrow\quad
        E=mc^2.
    \]
    \item Masseloses Teilchen:
    \[
        m=0
        \quad\Rightarrow\quad
        E=pc.
    \]
\end{itemize}
\end{conceptbox}

\begin{examplebox}
Fuer ein Photon gilt:
\[
    m=0.
\]
Damit folgt aus
\[
    E^2=p^2c^2
\]
direkt
\[
    E=pc.
\]
\end{examplebox}

\begin{examtipbox}
Diese Formel ist zentral:
\[
    E^2=p^2c^2+m^2c^4.
\]
Sie verbindet die messbaren Groessen Energie und Impuls mit der Masse des Teilchens.
\end{examtipbox}

\section{Natuerliche Einheiten}

In der Teilchenphysik verwendet man oft natuerliche Einheiten mit
\[
    c=1.
\]

\begin{conceptbox}
Wenn man
\[
    c=1
\]
setzt, werden Formeln kuerzer.

Dann wird aus
\[
    E^2=p^2c^2+m^2c^4
\]
einfach:
\[
    E^2=p^2+m^2.
\]
\end{conceptbox}

\begin{notebox}
In diesem Skript werden beide Schreibweisen verwendet.

Mit explizitem $c$:
\[
    E^2=p^2c^2+m^2c^4.
\]

In natuerlichen Einheiten:
\[
    E^2=p^2+m^2.
\]
\end{notebox}

\begin{examtipbox}
Wenn Massen in
\[
    \mathrm{MeV}/c^2
\]
und Impulse in
\[
    \mathrm{MeV}/c
\]
gegeben sind, ist die Schreibweise mit explizitem $c$ sinnvoll.

Wenn alle Groessen in
\[
    \mathrm{GeV}
\]
angegeben werden, wird haeufig stillschweigend $c=1$ verwendet.
\end{examtipbox}

\section{Viererimpuls}

In der speziellen Relativitaet fasst man Energie und Impuls zu einem Vierervektor
zusammen.

\begin{definitionbox}
Der \emph{Viererimpuls} ist
\[
    p^\mu
    =
    \left(
        \frac{E}{c},
        \vec p
    \right).
\]
\end{definitionbox}

\begin{conceptbox}
Der Viererimpuls kombiniert Energie und Impuls in einer Groesse, die sich unter
Lorentztransformationen einfach verhaelt.

Das ist besonders nuetzlich, weil Energie und Impuls in relativistischen Rechnungen
nicht getrennt betrachtet werden sollten.
\end{conceptbox}

\begin{formulabox}
Das Quadrat des Viererimpulses ist invariant:
\[
    p^\mu p_\mu
    =
    \frac{E^2}{c^2}
    -
    \vec p^{\,2}
    =
    m^2c^2.
\]
\end{formulabox}

\begin{conceptbox}
Invariant bedeutet:
Der Wert ist in allen Inertialsystemen gleich.

Energie und Impuls einzeln haengen vom Bezugssystem ab.
Die Kombination
\[
    E^2-p^2c^2
\]
ist aber invariant.
\end{conceptbox}

\section{Invariante Masse}

Die invariante Masse ist eine zentrale Groesse fuer Teilchenreaktionen.
Sie ist besonders wichtig, wenn ein kurzlebiges Teilchen zerfaellt und man es aus
seinen Zerfallsprodukten rekonstruieren will.

\begin{definitionbox}
Die \emph{invariante Masse} eines Systems ist die Masse, die sich aus der gesamten
Energie und dem gesamten Impuls des Systems ergibt.
\end{definitionbox}

\begin{formulabox}
Fuer ein System mit Gesamtenergie
\[
    E_{\mathrm{ges}}
\]
und Gesamtimpuls
\[
    \vec p_{\mathrm{ges}}
\]
gilt:
\[
    M^2c^4
    =
    E_{\mathrm{ges}}^2
    -
    p_{\mathrm{ges}}^2c^2.
\]
\end{formulabox}

\begin{formulabox}
In natuerlichen Einheiten gilt:
\[
    M^2
    =
    E_{\mathrm{ges}}^2
    -
    p_{\mathrm{ges}}^2.
\]
\end{formulabox}

\begin{conceptbox}
Die invariante Masse ist gleich in allen Inertialsystemen.
Deshalb kann man sie aus den gemessenen Energien und Impulsen der Endteilchen
berechnen.

Wenn ein Teilchen in mehrere Teilchen zerfaellt, ergibt die invariante Masse der
Zerfallsprodukte die Masse des urspruenglichen Teilchens.
\end{conceptbox}

\begin{examtipbox}
Typische Anwendung:
Ein kurzlebiges Teilchen wird nicht direkt gesehen.
Man misst seine Zerfallsprodukte und berechnet deren invariante Masse.

Ein Peak in der invarianten-Masse-Verteilung zeigt eine Resonanz oder ein neues
Teilchen.
\end{examtipbox}

\section{Invariante Masse von zwei Teilchen}

Betrachte zwei Teilchen mit Energien
\[
    E_1,\qquad E_2
\]
und Impulsen
\[
    \vec p_1,\qquad \vec p_2.
\]

\begin{formulabox}
Die invariante Masse des Zweiteilchensystems ist:
\[
    M^2c^4
    =
    (E_1+E_2)^2
    -
    \left|\vec p_1+\vec p_2\right|^2c^2.
\]
\end{formulabox}

\begin{conceptbox}
Wenn zwei Teilchen aus dem Zerfall eines Mutterteilchens stammen, dann ist
\[
    M
\]
die rekonstruierte Masse des Mutterteilchens.

Bei vielen Ereignissen erhaelt man eine Verteilung.
Eine Resonanz erscheint als Maximum in dieser Verteilung.
\end{conceptbox}

\begin{examplebox}
Bei der Suche nach einem Teilchen, das in ein Elektron-Positron-Paar zerfaellt,
misst man
\[
    e^-,
    \qquad
    e^+.
\]
Dann berechnet man:
\[
    M_{e^+e^-}^2c^4
    =
    (E_{e^+}+E_{e^-})^2
    -
    \left|\vec p_{e^+}+\vec p_{e^-}\right|^2c^2.
\]
Ein Peak in der Verteilung von
\[
    M_{e^+e^-}
\]
zeigt ein Teilchen, das in
\[
    e^+e^-
\]
zerfallen ist.
\end{examplebox}

\section{Schwerpunktssystem}

Das Schwerpunktssystem ist das Bezugssystem, in dem der Gesamtimpuls verschwindet.

\begin{definitionbox}
Das \emph{Schwerpunktssystem} ist das Bezugssystem, in dem der gesamte Impuls des
Systems null ist:
\[
    \vec p_{\mathrm{ges}}=0.
\]
\end{definitionbox}

\begin{conceptbox}
Im Schwerpunktssystem ist die gesamte Energie besonders einfach:
\[
    E_{\mathrm{CM}} = Mc^2.
\]

Hier ist $M$ die invariante Masse des gesamten Systems.
\end{conceptbox}

\begin{formulabox}
Im Schwerpunktssystem gilt:
\[
    E_{\mathrm{CM}}^2
    =
    M^2c^4.
\]
Allgemein gilt:
\[
    E_{\mathrm{CM}}^2
    =
    E_{\mathrm{ges}}^2
    -
    p_{\mathrm{ges}}^2c^2.
\]
\end{formulabox}

\begin{examtipbox}
Das Schwerpunktssystem ist wichtig, weil dort die Energie
\[
    E_{\mathrm{CM}}
\]
direkt die fuer neue Teilchen verfuegbare Energie beschreibt.
\end{examtipbox}

\section{Mandelstam-Variable s}

In der Teilchenphysik verwendet man haeufig die Variable $s$.

\begin{definitionbox}
Die Mandelstam-Variable $s$ ist das Quadrat der gesamten Schwerpunktsenergie in
natuerlichen Einheiten:
\[
    s = E_{\mathrm{CM}}^2.
\]
\end{definitionbox}

\begin{formulabox}
Mit Viererimpulsen gilt:
\[
    s
    =
    (p_1+p_2)^2.
\]
In natuerlichen Einheiten schreibt man:
\[
    E_{\mathrm{CM}}=\sqrt{s}.
\]
\end{formulabox}

\begin{conceptbox}
$s$ ist invariant.
Das bedeutet:
Man kann $s$ in jedem Bezugssystem berechnen und erhaelt denselben Wert.

Das ist praktisch, weil manche Rechnungen im Laborsystem einfacher sind, waehrend
die physikalische Interpretation im Schwerpunktssystem einfacher ist.
\end{conceptbox}

\section{Fixed-Target-Experiment}

Bei einem Fixed-Target-Experiment trifft ein bewegtes Teilchen auf ein ruhendes
Target.

\begin{definitionbox}
Ein \emph{Fixed-Target-Experiment} ist ein Experiment, bei dem ein beschleunigter
Strahl auf ein ruhendes Target trifft.
\end{definitionbox}

\begin{conceptbox}
Im Laborsystem hat das Target vor der Kollision keinen Impuls.
Ein grosser Teil der Energie des einfallenden Teilchens steckt nach der Kollision
in der Bewegung des Schwerpunktsystems.

Deshalb steht nicht die gesamte Strahlenergie fuer die Erzeugung neuer Teilchen zur
Verfuegung.
\end{conceptbox}

\begin{formulabox}
Fuer ein Projektil mit Gesamtenergie $E_1$ und Masse $m_1$, das auf ein ruhendes
Target der Masse $m_2$ trifft, gilt in natuerlichen Einheiten:
\[
    s
    =
    m_1^2
    +
    m_2^2
    +
    2m_2E_1.
\]
Damit:
\[
    E_{\mathrm{CM}}
    =
    \sqrt{s}.
\]
\end{formulabox}

\begin{examtipbox}
Bei Fixed-Target-Experimenten waechst die Schwerpunktsenergie fuer grosse
Strahlenergie nur ungefaehr wie:
\[
    E_{\mathrm{CM}}\propto \sqrt{E_1}.
\]

Das ist viel langsamer als beim Collider.
\end{examtipbox}

\section{Collider}

Bei einem Collider treffen zwei Teilchenstrahlen frontal aufeinander.

\begin{definitionbox}
Ein \emph{Collider} ist ein Beschleunigerexperiment, bei dem zwei gegenlaeufige
Teilchenstrahlen zur Kollision gebracht werden.
\end{definitionbox}

\begin{conceptbox}
Der Vorteil eines Colliders ist:
Der Gesamtimpuls der beiden Strahlen kann im Idealfall verschwinden.

Dann ist das Laborsystem gleichzeitig das Schwerpunktssystem.
Die Strahlenergie steht direkt als Schwerpunktsenergie zur Verfuegung.
\end{conceptbox}

\begin{formulabox}
Fuer zwei frontal kollidierende Teilchen gleicher Masse und gleicher Energie $E$
gilt bei hohen Energien naeherungsweise:
\[
    E_{\mathrm{CM}}
    \approx
    2E.
\]
\end{formulabox}

\begin{examplebox}
Beim LHC koennen zwei Protonenstrahlen mit je
\[
    E=\SI{7}{TeV}
\]
frontal kollidieren.

Dann ist die Schwerpunktsenergie:
\[
    E_{\mathrm{CM}}
    =
    \SI{14}{TeV}.
\]
\end{examplebox}

\begin{examtipbox}
Collider sind fuer Hochenergie-Teilchenphysik besonders effizient, weil die
Schwerpunktsenergie naeherungsweise linear mit der Strahlenergie waechst.
\end{examtipbox}

\section{Vergleich: Fixed Target und Collider}

\begin{center}
\begin{tabular}{|p{4.2cm}|p{5.3cm}|p{5.3cm}|}
\hline
Eigenschaft & Fixed Target & Collider \\
\hline
Geometrie & Strahl trifft ruhendes Target & Zwei Strahlen kollidieren frontal \\
\hline
Target & ruhend & zweiter Strahl \\
\hline
Schwerpunktsenergie & waechst langsam mit Strahlenergie & waechst naeherungsweise linear mit Strahlenergie \\
\hline
Vorteil & hohe Targetdichte moeglich & hohe Schwerpunktsenergie moeglich \\
\hline
Nachteil & viel Energie steckt in Schwerpunktsbewegung & technisch schwierige Strahlueberlappung \\
\hline
\end{tabular}
\end{center}

\begin{conceptbox}
Fixed-Target-Experimente sind oft gut fuer hohe Intensitaet und Praezisionsmessungen.

Collider sind besonders gut, wenn man moeglichst hohe Schwerpunktsenergie erreichen
will.
\end{conceptbox}

\begin{examtipbox}
Wenn nach dem Grund gefragt wird, warum Collider gebaut werden:

Ein Collider macht viel mehr der Strahlenergie im Schwerpunktssystem nutzbar.
Dadurch koennen schwerere Teilchen erzeugt werden.
\end{examtipbox}

\section{Schwellenenergie fuer Teilchenerzeugung}

Neue Teilchen koennen nur erzeugt werden, wenn genug Schwerpunktsenergie vorhanden ist.

\begin{definitionbox}
Die \emph{Schwellenenergie} ist die minimale Energie, die fuer eine Reaktion oder
Teilchenerzeugung erforderlich ist.
\end{definitionbox}

\begin{conceptbox}
Im Schwerpunktssystem ist die Schwelle besonders einfach:
Die Schwerpunktsenergie muss mindestens der Summe der Ruheenergien der Endteilchen
entsprechen.
\end{conceptbox}

\begin{formulabox}
Fuer eine Reaktion mit Endteilchenmassen
\[
    m_1,\quad m_2,\quad \ldots,\quad m_n
\]
gilt an der Schwelle:
\[
    E_{\mathrm{CM,thr}}
    =
    \left(
        m_1+m_2+\ldots+m_n
    \right)c^2.
\]
\end{formulabox}

\begin{conceptbox}
An der Schwelle haben die Endteilchen im Schwerpunktssystem keine zusaetzliche
kinetische Energie.
Die gesamte Energie steckt in ihren Ruhemassen.
\end{conceptbox}

\section{Resonanzen und invariante Masse}

Kurzlebige Teilchen erscheinen in Experimenten oft als Resonanzen.

\begin{definitionbox}
Eine \emph{Resonanz} ist ein kurzlebiger angeregter Zustand oder ein kurzlebiges
Teilchen, das in einer invarianten-Masse-Verteilung als Peak erscheint.
\end{definitionbox}

\begin{conceptbox}
Man misst die Endprodukte eines Zerfalls und berechnet ihre invariante Masse.
Wenn viele Ereignisse bei derselben invarianten Masse auftreten, spricht man von
einem Resonanzpeak.

Die Lage des Peaks gibt die Masse der Resonanz.
Die Breite des Peaks haengt mit der Lebensdauer zusammen.
\end{conceptbox}

\begin{formulabox}
Der Zusammenhang zwischen Breite $\Gamma$ und Lebensdauer $\tau$ ist grob:
\[
    \tau \sim \frac{\hbar}{\Gamma}.
\]
\end{formulabox}

\begin{examtipbox}
Breite Resonanz bedeutet kurze Lebensdauer.

Schmale Resonanz bedeutet lange Lebensdauer.
\end{examtipbox}

\section{Beispiel: Elektron-Positron-Annihilation}

Elektron-Positron-Collider sind besonders sauber, weil Elektron und Positron
punktfoermige Leptonen sind.

\begin{conceptbox}
Bei einer Elektron-Positron-Kollision kann die gesamte Schwerpunktsenergie in neue
Teilchen gehen.

Schematisch:
\[
    e^- + e^+ \rightarrow \gamma^* \rightarrow X.
\]

Bei geeigneter Schwerpunktsenergie kann auch ein schweres Boson oder eine Resonanz
erzeugt werden.
\end{conceptbox}

\begin{examplebox}
Wenn
\[
    E_{\mathrm{CM}}
\]
gleich der Ruheenergie eines Teilchens ist, kann dieses Teilchen resonant erzeugt
werden.

Zum Beispiel:
\[
    E_{\mathrm{CM}} \approx m_Zc^2
\]
fuer die Erzeugung eines $Z$-Bosons in
\[
    e^+e^-
\]
-Kollisionen.
\end{examplebox}

\section{Luminositaet}

Hohe Schwerpunktsenergie allein reicht nicht.
Man braucht auch genug Kollisionen.

\begin{definitionbox}
Die \emph{Luminositaet} $L$ beschreibt, wie intensiv zwei Teilchenstrahlen miteinander
kollidieren.
\end{definitionbox}

\begin{formulabox}
Die Ereignisrate $R$ eines Prozesses ist:
\[
    R = L\sigma.
\]
Dabei ist:
\begin{itemize}
    \item $R$ die Ereignisrate,
    \item $L$ die Luminositaet,
    \item $\sigma$ der Wirkungsquerschnitt des Prozesses.
\end{itemize}
\end{formulabox}

\begin{conceptbox}
Ein Prozess mit kleinem Wirkungsquerschnitt ist selten.
Um trotzdem genug Ereignisse zu messen, braucht man eine hohe Luminositaet.

Deshalb optimieren Collider nicht nur Energie, sondern auch:
\begin{itemize}
    \item Strahlintensitaet,
    \item Strahlfokussierung,
    \item Bunchstruktur,
    \item Kollisionsfrequenz.
\end{itemize}
\end{conceptbox}

\begin{examtipbox}
Energie entscheidet, ob ein Prozess kinematisch moeglich ist.

Luminositaet entscheidet, wie viele Ereignisse man pro Zeit beobachtet.
\end{examtipbox}

\section{Einheiten fuer Wirkungsquerschnitt und Luminositaet}

\begin{definitionbox}
Der Wirkungsquerschnitt hat die Dimension einer Flaeche.
Eine wichtige Einheit ist das Barn:
\[
    1\,\mathrm{barn}=10^{-28}\,\mathrm{m^2}.
\]
\end{definitionbox}

\begin{conceptbox}
In der Teilchenphysik sind Wirkungsquerschnitte oft sehr klein.
Daher verwendet man Untereinheiten:
\[
    \mathrm{mb},\quad \mu\mathrm{b},\quad \mathrm{nb},\quad \mathrm{pb},\quad \mathrm{fb}.
\]
\end{conceptbox}

\begin{formulabox}
Da
\[
    R=L\sigma
\]
die Einheit
\[
    \mathrm{s}^{-1}
\]
haben muss, hat die Luminositaet die Einheit:
\[
    \mathrm{cm}^{-2}\mathrm{s}^{-1}
\]
oder entsprechend
\[
    \mathrm{m}^{-2}\mathrm{s}^{-1}.
\]
\end{formulabox}

\section{Integrierte Luminositaet}

Neben der momentanen Luminositaet verwendet man die integrierte Luminositaet.

\begin{definitionbox}
Die \emph{integrierte Luminositaet} ist die ueber die Messzeit integrierte
Luminositaet:
\[
    \mathcal{L}_{\mathrm{int}}
    =
    \int L(t)\,\dd t.
\]
\end{definitionbox}

\begin{formulabox}
Die erwartete Anzahl von Ereignissen ist:
\[
    N
    =
    \mathcal{L}_{\mathrm{int}}\sigma.
\]
\end{formulabox}

\begin{conceptbox}
Die integrierte Luminositaet misst die insgesamt gesammelte Datenmenge eines
Experiments.

Je groesser
\[
    \mathcal{L}_{\mathrm{int}}
\]
ist, desto mehr seltene Prozesse kann man untersuchen.
\end{conceptbox}

\section{Impuls- und Energieerhaltung in Zerfaellen}

Relativistische Kinematik ist auch fuer Zerfaelle wichtig.

\begin{conceptbox}
Bei einem Zerfall muessen Energie und Impuls erhalten bleiben.
Im Ruhesystem des Mutterteilchens ist der Gesamtimpuls vor dem Zerfall null.

Daher muessen die Impulse der Zerfallsprodukte vektoriell zusammen null ergeben.
\end{conceptbox}

\begin{examplebox}
Bei einem Zweikoerperzerfall
\[
    A \rightarrow 1+2
\]
im Ruhesystem von $A$ gilt:
\[
    \vec p_1 + \vec p_2 = 0.
\]
Also:
\[
    \vec p_1 = -\vec p_2.
\]
Die Teilchen fliegen ruecken an ruecken auseinander.
\end{examplebox}

\begin{conceptbox}
Bei einem Dreikoerperzerfall verteilt sich die Energie variabel auf drei Teilchen.
Deshalb entstehen kontinuierliche Energiespektren.

Das ist wichtig fuer das Beta-Spektrum.
\end{conceptbox}

\section{Mini-Zusammenfassung}

\begin{notebox}
Die wichtigsten Punkte dieses Kapitels sind:
\begin{itemize}
    \item Der Lorentzfaktor ist
    \[
        \gamma=\frac{1}{\sqrt{1-\beta^2}}.
    \]
    \item Die Gesamtenergie ist
    \[
        E=\gamma mc^2.
    \]
    \item Der relativistische Impuls ist
    \[
        \vec p=\gamma m\vec v.
    \]
    \item Die Energie-Impuls-Beziehung ist
    \[
        E^2=p^2c^2+m^2c^4.
    \]
    \item Die invariante Masse eines Systems ist durch
    \[
        M^2c^4=E_{\mathrm{ges}}^2-p_{\mathrm{ges}}^2c^2
    \]
    gegeben.
    \item Im Schwerpunktssystem ist der Gesamtimpuls null.
    \item Die Schwerpunktsenergie ist die Energie, die fuer Teilchenerzeugung relevant ist.
    \item Fixed-Target-Experimente nutzen die Strahlenergie weniger effizient fuer Schwerpunktsenergie.
    \item Collider machen viel mehr Strahlenergie im Schwerpunktsystem verfuegbar.
    \item Luminositaet und Wirkungsquerschnitt bestimmen die Ereignisrate:
    \[
        R=L\sigma.
    \]
    \item Integrierte Luminositaet bestimmt die insgesamt erwartete Ereigniszahl:
    \[
        N=\mathcal{L}_{\mathrm{int}}\sigma.
    \]
\end{itemize}
\end{notebox}

\section{Pruefungscheck}

\begin{examtipbox}
Nach diesem Kapitel solltest du folgende Fragen beantworten koennen:
\begin{enumerate}
    \item Warum braucht man relativistische Kinematik in der Teilchenphysik?
    \item Was ist der Lorentzfaktor?
    \item Wie lautet die relativistische Gesamtenergie?
    \item Wie lautet der relativistische Impuls?
    \item Wie lautet die Energie-Impuls-Beziehung?
    \item Was gilt fuer masselose Teilchen?
    \item Was ist ein Viererimpuls?
    \item Was bedeutet invariant?
    \item Was ist die invariante Masse?
    \item Wie rekonstruiert man die Masse eines zerfallenden Teilchens?
    \item Was ist das Schwerpunktssystem?
    \item Was ist die Schwerpunktsenergie?
    \item Was ist die Mandelstam-Variable $s$?
    \item Was ist ein Fixed-Target-Experiment?
    \item Warum ist Fixed Target bei hohen Energien ineffizienter fuer Teilchenerzeugung?
    \item Was ist ein Collider?
    \item Warum erreicht ein Collider hoehere Schwerpunktsenergien?
    \item Was ist eine Schwellenenergie?
    \item Was ist eine Resonanz?
    \item Wie haengt Resonanzbreite mit Lebensdauer zusammen?
    \item Was ist Luminositaet?
    \item Wie haengt die Ereignisrate mit Luminositaet und Wirkungsquerschnitt zusammen?
    \item Was ist integrierte Luminositaet?
\end{enumerate}
\end{examtipbox}